
%\vspace{-1cm}
\setcounter{section}{-1}
\section{Introduction}
L'analyse numérique est un vaste champ d'étude des mathématiques qui a pour objectif de fournir des solutions approchées à des problèmes concrets. On cherche à mettre en place des méthodes numériques -souvent algorithmiques- qui permettent d'approximer le résultat de certains problèmes.%, parmi lesquels ont peut citer l'évaluation d'une fonction en un point, la résolution d'équations et de système d'équations, le calcul d'une intégrale\footnote{Le concept \guillemets{d'intégrales} sera présenté en cours de mathématique de 4\up{ème} année.}, l'interpolation\footnote{\guillemets{L'interpolation} tente de retrouver l'expression algébrique d'une fonction passant par des points donnés.} ou encore l'optimisation.


\medskip Les équations sont des objets que l'on rencontre fréquemment en mathématiques et l'intérêt pour leur résolution apparaît dans de nombreux contextes scientifiques. Dans la passé, la résolution numérique d'équations a été un moteur pour le développement de plusieurs notions  mathématiques et c'est probablement un des domaines les plus anciens de l'analyse numérique.

\medskip La plupart des phénomènes physiques, chimiques ou biologiques issus de la technologie moderne sont régis par des équations complexes. C'est en étudiant et en observant la nature de ces phénomènes que nous parvenons à modéliser leur fonctionnement à l'aide d'outils mathématiques. Cependant, la recherche de solutions à ces équations peut être (et est très souvent!) compliquée car la modélisation fait intervenir beaucoup de paramètres.


\medskip Le développement, relativement récent, de puissants moyens de calculs a ouvert de nouvelles pistes d'explorations ainsi que de nouvelles perspectives de recherches pour répondre aux besoins de l'ingénierie. Cette dernière a besoin de méthodes efficaces pour résoudre des équations et  la résolution numérique de ces équations à l'aide d'un ordinateur nécessite des connaissances approfondies en mathématiques.

\medskip
Nous allons étudier trois méthodes permettant de résoudre numériquement des équations.

%Nous le travaillons activement en mathématiques mais la recherche d'une quantité manquante est un questions récurrente dans  toutes les autres sciences (au sens large).

%\medskip Nombre de phénomènes physiques, chimiques, économiques, sociaux, démographiques, naturels, et la liste est encore longue, sont modélisés à l'aide d'équations qui permettent de décrire non seulement la nature du phénomène mais également de répondre à certaines questions et résoudre des problèmes.  

